\section{Introduction}\label{sec:introduction}

Multiple ergodic averages measure how several observations along one orbit
interact.  We consider the three observations at times $n,2n,3n$.  Let
$(X,\mathcal F,\mu)$ be a probability space and let $T:X\to X$ be an
invertible measure preserving transformation such that both $T$ and $T^{-1}$
are measurable.  For bounded measurable functions $f_1,f_2,f_3:X\to\C$, put
\begin{equation}
 A_N(f_1,f_2,f_3)(x)
 =\frac1N\sum_{n=1}^N
       f_1(T^n x)f_2(T^{2n}x)f_3(T^{3n}x),\qquad N\ge1.
 \label{eq:main-average}
\end{equation}
The transformation is \emph{mixing} if, for every $A,B\in\mathcal F$,
\begin{equation}
 \mu(A\cap T^{-m}B)\longrightarrow\mu(A)\mu(B)
 \quad\text{as }|m|\to\infty,\quad m\in\Z.
 \label{eq:mixing}
\end{equation}

\begin{theorem}\label{thm:main}
Let $(X,\mathcal F,\mu)$ be an arbitrary probability space and let
$T:X\to X$ be an invertible mixing probability preserving transformation
with measurable inverse.  For every triple of bounded measurable functions
$f_1,f_2,f_3:X\to\C$,
\[
 A_N(f_1,f_2,f_3)(x)
 \longrightarrow \prod_{j=1}^3\int_X f_j\,d\mu
 \quad\text{for $\mu$-almost every }x
\]
as $N\to\infty$ through all positive integers.  The exceptional null set
may depend on the triple of functions.
\end{theorem}

The analytic input to the proof is the companion paper
\emph{An $L^3$ bound for the trilinear Hilbert transform} by OpenAI
\cite{H}, denoted below by H.  We use its local-form construction and
several estimates within its proof.  The required estimates are stated at
their points of use.  We also use H's uniform local bound for fixed compactly supported
scalar $L^q$ inputs at one exponent $q\in(2,3)$ shared by its argument and ours.
Section~\ref{sec:completion} justifies this common choice using the
parameter order in H's proof.  The new work here is an extension that allows
the inputs to vary with spatial depth.  Theorem~\ref{thm:depth-local} states
that extension, and
Sections~\ref{sec:rows}--\ref{sec:completion} prove it from the specified
results of H.

Furstenberg's ergodic proof of Szemer\'edi's theorem placed arithmetic
progressions of orbit times at the center of ergodic Ramsey theory
\cite[Theorem~1.4]{Furstenberg1977}.  The convergence theory subsequently
developed along two distinct lines.  For norm convergence, Conze and Lesigne
proved the three-function result under total ergodicity
\cite{ConzeLesigne1987}, and Host and Kra removed that restriction
\cite{HostKra2001}.  Host and Kra then proved norm convergence for every
finite number of functions using nilpotent characteristic factors
\cite[Theorem~1.1]{HostKra2005}; Ziegler gave an independent, later
construction and an alternative proof \cite{Ziegler2007}.  For weakly mixing
transformations the product-valued norm limit was already classical.  We
include the van der Corput induction for that fact, specialized to the
mixing hypothesis \eqref{eq:mixing}; see
\cite[Section~3]{FKO1982} and \cite[Section~3]{Bergelson1987}.

Pointwise convergence asks for control along almost every individual orbit.
Bourgain proved the two-function theorem for two powers of one transformation
\cite{Bourgain1990}.  Higher-function pointwise results have been established
under additional structure, including $K$ systems on Lebesgue probability spaces
\cite[Theorem~4.2]{DerrienLesigne1996}, a singular-spectrum condition on the
Pinsker factor of a weakly mixing system \cite{Assani1998}, and weak mixing
together with the PID property, which requires every pairwise-independent
finite self-joining to be the product \cite{GHSY2018}.
The last result includes finite-rank mixing transformations
\cite{GHSY2018}.  Ergodic measurable distal systems on Lebesgue probability
spaces provide another class with pointwise convergence \cite{HSY2019}.
The September 2026 account of Kosz, Mirek, Peluse, Wan, and Wright
discusses the unrestricted three-function problem for the times $n,2n,3n$
\cite[Section~1.8, item~5]{KMPWW2026}; their theorem treats polynomials of
distinct degrees.  Theorem~\ref{thm:main} concerns the mixing class.

The harmonic-analysis setting has a related progression structure.  In the
four-linear form associated with the trilinear Hilbert transform, matched
quadratic phases can persist because
\[
 \sum_{j=0}^3 \kappa_j(x+jt)^a=0\quad(a=0,1,2),
 \qquad (\kappa_0,\kappa_1,\kappa_2,\kappa_3)=(-1,3,-3,1).
\]
The bilinear time-frequency theory of Lacey and Thiele
\cite{LT1997,LT1999} is a predecessor.  Tao connected cancellation for
truncated multilinear Hilbert transforms with higher-order uniformity
\cite{Tao2015}; Zorin-Kranich and Durcik, Kova\v c, and Thiele developed
further cancellation estimates through simplex forms
\cite{ZK2017,DKT2019}.  H uses the quantitative degree-two inverse theorem
of Leng, Sah, and Sawhney \cite[Theorem~1.2]{LSS2024} in a construction that
sums local four-term progression forms uniformly over their depths.  Our
argument extends that local construction to the inputs needed for oscillation.

\subsection{Proof strategy}

By homogeneity, assume $|f_j|\leq1$.  For an averaging profile $\varphi$
(Definition~\ref{def:averaging-profile}), which is a compactly supported
smooth real function of integral one, write
$\varphi_a(t)=a^{-1}\varphi(t/a)$ and consider the finite sums
\[
 V_k^\varphi(x)=\sum_{n\in\Z}\varphi_{2^k}(n)
           \prod_{j=1}^3 f_j(T^{jn}x),\qquad k\ge0.
\]
Proposition~\ref{prop:orbit-oscillation} proves that, for every integer
$W\geq1$, every list $0\le k_1<\cdots<k_{W+1}$ fixed independently of
$x$ satisfies
\begin{equation}
 \sum_{i=1}^W\int_X
    \max_{k_i\le\ell\le k_{i+1}}
        |V_\ell^\varphi-V_{k_i}^\varphi|\,d\mu
 \le C_\varphi\sqrt W.
 \label{eq:intro-oscillation}
\end{equation}
If the averages were not Cauchy on a set of positive measure, fixed endpoints
could be chosen so that every block contributes a uniformly positive integral
oscillation.  The left side would then grow proportionally to $W$, contradicting
\eqref{eq:intro-oscillation}.  This is the quantitative oscillation method
of Bourgain \cite[Section~2C]{Bourgain1989}, also used for bilinear singular
integrals by Demeter \cite{Demeter2007}.  For each fixed profile, it gives
almost-everywhere convergence of $V_k^\varphi$ as $k\to\infty$ on every
invertible probability preserving system.  Under \eqref{eq:mixing}, the norm
argument and summation by parts for smooth weights identify this limit as
$p=\prod_j\int_X f_j\,d\mu$, independently of the profile.  Signed $L^1$
approximations to interval weights and finite rational grids then recover
all positive integer lengths in \eqref{eq:main-average}.

The pointwise maximum creates the new analytic requirement.  At each point
$x$, choose $\ell_i(x)\in\{k_i,\ldots,k_{i+1}\}$ attaining the $i$th maximum.
The identity
\[
 V_{\ell_i(x)}^\varphi(x)-V_{k_i}^\varphi(x)
 =\sum_{k=k_i}^{\ell_i(x)-1}
       \bigl(V_{k+1}^\varphi(x)-V_k^\varphi(x)\bigr)
 =:d_i(x)
\]
shows that the selected difference is a sum over a prefix of the scale block.
Choose $|\alpha_i(x)|\leq1$ with $\alpha_i(x)d_i(x)=|d_i(x)|$.  The normalized
test at scale $k$ is
\[
 \frac{\alpha_i(x)}{4\sqrt W}\,
       \mathbf1_{\{k_i\leq k<\ell_i(x)\}}.
\]
Its supremum plus internal total variation on the block is at most
$1/(2\sqrt W)$, since its size is at most $1/(4\sqrt W)$ and it has at most
one jump.  In the real-variable model, H's local forms integrate
$\prod_{j=0}^3z_j(y+jt)$ against smooth kernels supported where all four
coordinates lie inside one dyadic interval.
In the application only the zeroth input varies with depth; the other three
inputs use the same functions at every depth.  Theorem~\ref{thm:depth-local}
gives a local bound independent of the depth range and $W$, provided the
varying input has supremum plus variation at most $W^{-1/2}$ on each of its
$W$ fixed blocks.  Reversal from increasing scale order to spatial depth
preserves this quantity.  Section~\ref{sec:transfer} realizes the scale
differences by normalized translation averaging of local kernels, then passes
through integer sequences to finite orbit segments.  Multiplying the normalized test back by
$4\sqrt W$ gives \eqref{eq:intro-oscillation}.

Section~\ref{sec:rows} supplies the operator step.  The structured operators
$A_\nu$ arising in H, called rows, have a common $L^2$ domain and obey the
Bessel bound
\[
 \sum_\nu\|A_\nu f\|^2\leq B^2\|f\|_2^2
 \qquad\text{for every }f\in L^2,
\]
with the same input in every term.  Lemma~\ref{lem:row-maximal} bounds the
$L^1$ norm of the pointwise maximum of tested adjoint partial sums in terms
of the tests' $\ell^2$ size and the Bessel bound, with logarithmic dependence
on the column budget and no dependence on the number of depths.  The tested
adjoints are, on each relevant cell, combinations with spatially constant
coefficients of column functions retained on descendant cells.  Their distinct
orthogonal auxiliary coordinates control the coefficient vectors, and
refinement of the spatial cells lets the stopping argument assign each row to
one group on its entire support.  When each round reads one
depth and those depths are nondecreasing, pointwise summation by parts pairs
the partial sums with the input's size and internal variation on each block.
The $L^1$ estimate and Cauchy--Schwarz over blocks use the $W^{-1/2}$
normalization to remove dependence on $W$, yielding Lemma~\ref{lem:row-depth}.
Lemma~\ref{lem:smooth-depth} treats the smooth rows.
Sections~\ref{sec:base}--\ref{sec:completion} insert these estimates at the
specified steps of H that compare inputs at different depths and complete
the local estimate.
